% =======================================================================
% TECHNICAL APPENDIX -- Active Causal Experimentalist (ACE)
% NeurIPS 2026. No length limit on appendix.
% =======================================================================

% =======================================================================
\section{Notation and Problem Formulation (Extended)}
\label{app:notation}
% =======================================================================

We adopt Pearl's causal framework~\citep{pearl2009causality,pearl1995causal}. A Structural Causal Model (SCM) $\mathcal{M} = \langle \mathcal{U}, \mathcal{V}, \mathcal{F}, P(\mathcal{U}) \rangle$ consists of exogenous variables $\mathcal{U} = \{U_1, \ldots, U_m\}$, endogenous variables $\mathcal{V} = \{V_1, \ldots, V_n\}$, structural equations $\mathcal{F} = \{f_1, \ldots, f_n\}$ where $V_i = f_i(\text{Pa}_i, U_i)$, and a distribution $P(\mathcal{U})$.

The causal relationships induce a directed acyclic graph (DAG) $\mathcal{G} = (\mathcal{V}, \mathcal{E})$ with $(V_j, V_i) \in \mathcal{E}$ iff $V_j \in \text{Pa}_i$. The observational distribution factorizes as
\begin{equation}
P(V_1, \ldots, V_n) = \prod_{i=1}^n P(V_i \mid \text{Pa}_i).
\end{equation}
An intervention $\text{do}(V_i = \nu)$ replaces $f_i$ with the constant $\nu$ and removes all edges into $V_i$ in the mutilated graph $\mathcal{G}_{\overline{V_i}}$ \citep{pearl2009causality}. This is a Rung-2 operation in Pearl's causal hierarchy: it requires active manipulation of the system rather than mere observation, and yields samples from the post-intervention distribution $P(V_j \mid \text{do}(V_i = \nu))$ for all $j \ne i$.

The impossibility of mechanism identification from observational data alone follows from Reichenbach's common-cause principle \citep{reichenbach1956}: for any bivariate association $A \perp\!\!\!\!\not\perp B$, there exist at least three causal structures ($A \to B$, $B \to A$, $A \leftarrow C \rightarrow B$) that are Markov equivalent \citep{spirtes2000causation,peters2017elements}. Distinguishing structural equations $f_i$ requires interventional data \citep{peters2017elements}.

\textbf{ACE's task.} Given known $\mathcal{G}$, ACE learns the structural equations $\{f_i\}$ by sequentially selecting interventions $c_t = \text{do}(V_{i_t} = \nu_t)$ to minimize the total mechanism estimation loss. The policy $\pi_\phi(c_t \mid s_t)$ adapts based on observed losses $\{L_i\} = \{L_1, \ldots, L_n\}$ where $L_i$ measures current discrepancy between the student SCM mechanism $\hat{f}_i$ and the true $f_i$.

% =======================================================================
\section{Implementation Details}
\label{app:impl}
% =======================================================================

\subsection{Ground Truth SCMs}
\label{app:scms}

\textbf{5-node benchmark.} The benchmark used throughout the main paper has structural equations
\begin{align*}
X_1 &\sim \mathcal{N}(0, 1) & \text{(exogenous root)} \\
X_4 &\sim \mathcal{N}(2, 1) & \text{(exogenous root)} \\
X_2 &= 2 X_1 + 1 + \varepsilon_2 & \text{(linear)} \\
X_3 &= 0.5 X_1 - X_2 + \sin(X_2) + \varepsilon_3 & \text{(nonlinear collider)} \\
X_5 &= 0.2 X_4^2 + \varepsilon_5 & \text{(quadratic)}
\end{align*}
with $\varepsilon_i \sim \mathcal{N}(0, 0.01)$. The graph has two roots ($X_1, X_4$), two intermediate nodes ($X_2, X_5$), and one collider ($X_3$ with parents $X_1, X_2$). The collider tests the policy's ability to identify the nonlinear interaction term $\sin(X_2)$, which requires observing $X_3$ under targeted interventions on both $X_1$ and $X_2$. The disconnected chain $X_4 \to X_5$ with a quadratic mechanism tests mechanism estimation for polynomial functions.

\textbf{30-node hierarchical SCM.} Constructed at runtime with seed-controlled coefficients in $[0.3, 0.7]$. Five-layer structure (see Figure~\ref{fig:30node-scm}):
\begin{itemize}[leftmargin=1.2em,topsep=0.2em,itemsep=0.1em]
\item Layer 0: 5 root nodes, $X_i \sim \mathcal{N}(0,1)$.
\item Layer 1 (X6--X10): 1--2 root parents each, linear combination with noise std $0.15$.
\item Layer 2 (X11--X20): Mix of single-parent and two-parent (collider) nodes; every third node is a collider.
\item Layer 3 (X21--X25): Complex colliders with 2--3 parents from Layer 2; sine nonlinearity on every fifth node.
\item Layer 4 (X26--X30): Single-parent leaf nodes.
\end{itemize}
The seed controls the random graph wiring within these constraints, ensuring reproducibility.

\textbf{Duffing oscillators.} Three coupled Duffing oscillators with true coupling $X_1 \leftrightarrow X_2 \leftrightarrow X_3$. The synchronization-induced correlation between $X_1$ and $X_3$ is a spurious association (Rung~1) with no direct causal edge; interventions on $X_2$ are necessary to break this correlation and estimate the direct couplings.

\subsection{Learner Architecture}
\label{app:learner}

The student SCM uses a \texttt{ModuleDict} of per-node mechanisms:
\begin{itemize}[leftmargin=1.2em,topsep=0.2em,itemsep=0.1em]
\item \textbf{Root nodes}: Learnable parameters $(\hat{\mu}, \hat{\sigma})$; loss is KL divergence between $\mathcal{N}(\hat{\mu}, \hat{\sigma}^2)$ and empirical samples.
\item \textbf{Non-root nodes}: 2-layer MLP with hidden dim 64 and ReLU activations, taking the parent vector $\text{Pa}_i$ as input and outputting a scalar prediction of $V_i$. Loss is MSE against observed $V_i$ under the current intervention.
\end{itemize}
Optimization: Adam~\citep{kingma2015adam}, learning rate $2 \times 10^{-3}$, batch size 100. Per-step training runs 100 epochs on the latest interventional batch; observational training every 3 steps runs 100 epochs on 200 fresh observational samples to maintain root-distribution estimates.

\subsection{Policy Architecture and Prompt Design}
\label{app:policy}

The policy $\pi_\phi$ is Qwen2.5-1.5B \citep{qwen2_5}, accessed via the HuggingFace Transformers API. Sampling: temperature 0.7, max new tokens 64. The text prompt has the form:
\begin{quote}\small
\texttt{Graph: X1->X2, X1->X3, X2->X3, X4->X5}\\
\texttt{Per-node losses: X1=0.92, X2=0.01, X3=1.45, X4=0.98, X5=0.07}\\
\texttt{Recent interventions (last 5): DO X2=1.3 (gain=0.5), ...}\\
\texttt{Propose your next intervention as: DO Xi = v}
\end{quote}
We generate $K = 4$ candidates per step by re-sampling at temperature 0.7. The prompt deliberately uses the same node labels as the SCM, allowing the LM's pretrained world-model prior to activate when node names are semantically meaningful (e.g., economic or genetic variables).

\subsection{Reward Components in Full}
\label{app:reward}

The reward in Equation~\eqref{eq:reward} has three components. Let $V_i$ be the targeted node, $\{L_j\}$ the per-node losses, and $H$ the recent-intervention history.

\textbf{Information gain.} $\Delta\mathcal{L}(c) = L_\text{before} - L_\text{after}$ where $L = \sum_i L_i$ is the validation loss summed across nodes, estimated by lookahead on a cloned learner.

\textbf{Node importance.} $w(V_i, \{L_j\}) = \texttt{cov\_bonus} \cdot \tfrac{L_i}{\sum_j L_j} \cdot (1 + n_i(H))^{-1/2}$ where $n_i(H)$ is the count of recent interventions on $V_i$. This combines a per-node weight (proportional to current loss, i.e., how poorly understood the mechanism is) with an under-sampling correction that prevents over-concentration on a single node.

\textbf{Diversity.} $D(V_i, H)$ rewards interventions on under-sampled nodes and under-explored intervention values; full definition in the public code repository.

Default weights: $\alpha = 0.1$ (\texttt{cov\_bonus=60.0}), $\gamma = 0.05$. These keep $\Delta\mathcal{L}$ dominant ($\sim$80--90\% of total reward) while ensuring sufficient exploration pressure.

\subsection{Bayesian OED Baseline Implementation}
\label{app:boed-impl}

Our Bayesian OED baseline approximates the expected information gain over the same parametric mechanism class as ACE's student. At each step:
\begin{enumerate}[leftmargin=1.5em]
\item For each candidate $(V_i, \nu)$ with $\nu$ uniformly sampled from $[-5, 5]$:
  \begin{enumerate}
    \item Sample $M=10$ posterior parameter draws $\theta^{(m)} \sim P(\theta \mid \mathcal{D}_\text{seen})$ by training $M$ student copies from different random initializations.
    \item For each draw, simulate outcome $y^{(m)} \sim P(\cdot \mid \theta^{(m)}, \text{do}(V_i = \nu))$.
    \item Estimate $\widehat{\Delta\mathcal{L}} = \frac{1}{M}\sum_m [L(\theta^{(m)} \mid \mathcal{D}_\text{seen}) - L(\theta^{(m)} \mid \mathcal{D}_\text{seen} \cup \{(V_i, \nu, y^{(m)})\})]$.
  \end{enumerate}
\item Execute the candidate with highest $\widehat{\Delta\mathcal{L}}$.
\end{enumerate}
We use $|\mathcal{V}| \times 20$ candidates per step. This faithfully adapts the ABCI \citep{toth2022active} and CBED \citep{tigas2022interventions} acquisition-function framework to the mechanism estimation setting where graph structure is known. Each Bayesian OED step takes $\sim$4--6 minutes on a CPU, limiting us to N=3 seeds.

% =======================================================================
\section{Theoretical Analysis (Extended)}
\label{app:theory}
% =======================================================================

\subsection{Full Proof of Proposition~\ref{prop:scale-invariance}}
\label{app:proof}

Recall Assumption~\ref{ass:diminish}: $r_t(a) = f(t) \cdot g(a)$ where $f(t) > 0$ is monotonically non-increasing.

\textbf{Part 1 (preference invariance).} For any $a_w, a_l$ with $g(a_w) > g(a_l)$:
\[r_t(a_w) - r_t(a_l) = f(t)(g(a_w) - g(a_l)) > 0 \quad \forall t,\]
since $f(t) > 0$. The preference ordering is therefore constant across all steps.

\textbf{Part 2 (value non-stationarity).} Under standard MDP value definitions,
\[V_t(s) = \mathbb{E}_a[r_t(a) \mid s] = f(t) \cdot \mathbb{E}_a[g(a) \mid s].\]
With $f$ decreasing, $V_t(s)$ is non-stationary. A value function approximator $\hat{V}(s)$ trained at episode $t$ misestimates $V$ at episode $t'$ by factor $f(t')/f(t)$. Empirically, $f(t')/f(t) \approx 1/500$ between early and late training, so the critic accumulates large value-estimation error even with frequent updates.

\textbf{Part 3 (DPO invariance).} The DPO loss is
\[\mathcal{L}_\text{DPO} = -\mathbb{E}_{(y_w, y_l)} \log \sigma\!\left(\beta \left[\log\tfrac{\pi_\phi(y_w)}{\pi_\text{ref}(y_w)} - \log\tfrac{\pi_\phi(y_l)}{\pi_\text{ref}(y_l)}\right]\right).\]
The argument depends only on policy log-ratio differences. Both $\pi_\phi$ and $\pi_\text{ref}$ are normalized distributions; their log-ratios depend on action quality $g(a)$ but not on the reward scale $f(t)$. Thus the gradient $\nabla_\phi \mathcal{L}_\text{DPO}$ is invariant to any multiplicative rescaling of all rewards by a positive constant. \hfill$\square$

\subsection{Convergence Rate Equivalence under Stationary $g$}
\label{app:convergence}

Under Assumption~\ref{ass:diminish} and standard DPO convergence assumptions \citep{rafailov2023direct}, ACE inherits the regret bound of preference-based bandits: after $T$ pairwise comparisons,
\[\mathbb{E}\!\left[\sum_{t=1}^T (g(a_\text{opt}) - g(a_t))\right] = \tilde{O}(\sqrt{T \log |\mathcal{A}|}),\]
matching the minimax regret for preference-based learning. Value-based RL achieves the same bound only when the value approximator tracks the non-stationary scaling $f(t) \cdot g(a)$; in our regime (where $f(t)$ varies by $500\times$ over the training horizon), this fails unless the critic update rate exceeds the rate of change of $f$.

\subsection{Connection to Expected Information Gain}
\label{app:info-gain}

The information gain term $\Delta\mathcal{L}(c)$ is a Monte Carlo estimate of the mutual information between the posterior over mechanism parameters $\theta$ and the data $\mathcal{D}_c$ generated by intervention $c$:
\begin{align}
I(\theta; \mathcal{D}_c \mid \sigma) &= H(\theta \mid \sigma) - H(\theta \mid \mathcal{D}_c, \sigma) \\
&\approx L(\theta \mid \sigma) - L(\theta \mid \sigma \cup \mathcal{D}_c),
\end{align}
where the approximation uses the local-quadratic (Laplace) relationship between negative log-likelihood loss $L$ and posterior entropy $H$ \citep{rainforth2024modern}. With Gaussian noise mechanisms and a flat prior, this approximation is tight near the posterior mode. The node importance term $w$ acts as a non-uniform prior weighting, upweighting mechanisms with high remaining uncertainty, approximating posterior-entropy-reduction-weighted EIG. ACE thus pursues the same EIG objective as Bayesian methods through policy learning, trading exact inference for scalability.

% =======================================================================
\section{Extended Experimental Results}
\label{app:experiments}
% =======================================================================

\subsection{5-Node Benchmark: Per-Seed Results}
\label{app:5node-detail}

Table~\ref{tab:5node-perseed} reports per-seed total losses for heuristic baselines and PPO (N=5 CURC seeds: 141, 271, 314, 577, 618) and Bayesian OED (N=3 seeds: 42, 123, 456). ACE per-seed data comes from the original production runs (seeds 42, 123, 456, 789, 1011). All runs use 171 episodes.

\begin{table}[h]
\centering
\caption{Per-seed total loss on the 5-node benchmark (171 episodes). Heuristic baseline and PPO columns use seeds 141/271/314/577/618 (CURC). Bayesian OED uses seeds 42/123/456. ACE uses seeds 42/123/456/789/1011 (original Lambda Cloud runs; a rerun on CURC was initiated but OOM'd on L40 GPUs due to $3$GB policy checkpoints).}
\label{tab:5node-perseed}
\small
\begin{tabular}{@{}l ccc cc@{}}
\toprule
Seed & Random & Round-Robin & Max-Var.\ & Bayesian OED & PPO \\
\midrule
141 & 2.090 & 2.186 & 2.280 & ---   & 2.062 \\
271 & 2.240 & 2.023 & 2.050 & ---   & 1.994 \\
314 & 2.227 & 2.238 & 2.104 & ---   & 2.133 \\
577 & 2.091 & 1.978 & 2.184 & ---   & 2.134 \\
618 & 2.210 & 2.088 & 1.917 & ---   & 2.062 \\
42  & ---   & ---   & ---   & 1.995 & ---   \\
123 & ---   & ---   & ---   & 2.181 & ---   \\
456 & ---   & ---   & ---   & 1.957 & ---   \\
\midrule
Mean & 2.171 & 2.103 & 2.107 & 2.044 & 2.077 \\
Std  & 0.075 & 0.107 & 0.131 & 0.120 & 0.057 \\
\bottomrule
\end{tabular}
\end{table}

The Bayesian OED mean of $2.044$ lies between Round-Robin ($2.103$) and Max-Variance ($2.107$), confirming it is the strongest passive/principled baseline. Its advantage over random baselines is modest ($\sim$6\%) compared to ACE's advantage ($\sim$70\%), which arises from multi-step learned strategy rather than single-step acquisition.

\subsection{30-Node SCM: Per-Seed Results}
\label{app:30node-detail}

Table~\ref{tab:30node-perseed} reports all per-seed results for the 30-node benchmark, including the locally-run Round-Robin and Max-Variance CPU baselines.

\begin{table}[h]
\centering
\caption{Per-seed results on the 30-node hierarchical SCM. ACE and Random both use the MLP-based student learner from the main ACE infrastructure (same evaluation protocol). ACE reports best-loss-achieved across training; Random reports final-episode loss. Round-Robin and Max-Variance require GPU execution with the same MLP learner for a valid comparison and are not reported here.}
\label{tab:30node-perseed}
\small
\begin{tabular}{@{}l cc@{}}
\toprule
Seed & ACE (best-loss, MLP learner) & Random (final, MLP learner) \\
\midrule
42   & 2.785 & 11.808 \\
123  & 1.798 & 12.225 \\
456  & 1.274 & 9.625  \\
789  & ---   & 11.041 \\
1011 & ---   & 11.729 \\
\midrule
Mean $\pm$ Std & $1.952\pm0.767$ (N=3) & $11.286\pm1.021$ (N=5) \\
\bottomrule
\end{tabular}
\smallskip\\
{\footnotesize Both methods use the same per-node MLP mechanism estimator, making the comparison fair. A CPU proxy run with a running-mean predictor gave different absolute numbers (Random: $14.70\pm7.15$) due to the different learner; those numbers are not comparable to ACE and are excluded.}
\end{table}

Seed 42 ACE exhibited training instability (loss oscillating $2.7$--$3.5$ throughout the 40 completed episodes); seeds 123 and 456 converged cleanly to $1.80$ and $1.27$. The stable-seed mean is $1.54$; we report the conservative N=3 mean of $1.95$ in the main text.

\textbf{Why Round-Robin and Max-Variance are not reported on 30-node.} These baselines require the same MLP mechanism learner as ACE for a valid comparison (the learner defines the loss function that is being minimized). Running 300 episodes $\times$ 25 steps/episode with a 30-node MLP learner requires $\sim$4--6 h per seed on CPU, or $\sim$2 h per seed on GPU. We ran them on CPU with a simple running-mean predictor (script: \texttt{scripts/runners/run\_30node\_cpu\_baselines.py}) but the results are not comparable to ACE and are therefore not included in the main analysis. Adding proper 30-node baseline comparisons with the MLP learner is the second highest-priority follow-up experiment (after graph misspecification).

\subsection{Component Ablation Detail}
\label{app:ablation-detail}

The component ablation (Table~\ref{tab:ablations}) was run as a 3-seed pilot (seeds 42, 123, 456) due to GPU compute constraints; the No-diversity ablation was run on N=5 seeds during an earlier reward-shaping investigation. Below we give qualitative failure-mode analysis for each configuration.

\textbf{No DPO} replaces the LM policy with random proposal generation while keeping the lookahead and select-best mechanism. The result ($2.10\pm0.11$) matches Round-Robin, confirming that the ability to \emph{evaluate} candidates (via lookahead) provides no benefit without \emph{learned} proposal generation. Random candidates---even when curated by lookahead---do not consistently include strategically important ones (e.g., concentrating on collider parents).

\textbf{No per-node convergence} forces termination at a fixed 100 episodes. The collider mechanism $X_3$ is typically not learned at episode 100; $L_{X_3}$ remains elevated and dominates total loss.

\textbf{No dedicated root learner} disables the observational training loop for root nodes. Their losses ($L_{X_1}, L_{X_4}$) remain elevated because direct interventions on roots replace $X_i$ with a constant, providing no signal about the root's natural distribution $\mathcal{N}(\mu, \sigma^2)$.

\textbf{No diversity} ($\gamma=0$) causes the policy to collapse onto a single high-information node. The total loss worsens beyond all baselines ($2.82$) because the un-targeted nodes are never learned, even if the targeted node is learned well.

\subsection{Duffing Oscillators: Full Results}
\label{app:duffing}

\begin{table}[h]
\centering
\caption{Coupled Duffing oscillator coupling-parameter estimation error (N=5 seeds, 100 episodes). Lower is better. The $X_1 \leftrightarrow X_3$ spurious correlation is an association-level artefact (Rung~1) without a direct causal edge; it is broken by clamping $X_2$ (Rung~2 intervention).}
\label{tab:duffing-full}
\small
\begin{tabular}{@{}lcccc@{}}
\toprule
Method & Coupling Error & Std & 95\% CI & Improvement vs.\ Random \\
\midrule
\textbf{ACE} & \textbf{0.042} & 0.036 & [0.006, 0.078] & $5.8\times$ \\
Round-Robin  & 0.238 & 0.076 & [0.162, 0.314] & $1.0\times$ \\
Random       & 0.245 & 0.121 & [0.124, 0.366] & --- \\
\bottomrule
\end{tabular}
\end{table}

ACE concentrates $62\%$ of interventions on $X_2$ (the intermediate oscillator) compared to $33\%$ under uniform allocation. Clamping $X_2$ decouples the synchronized chain, breaking the Rung-1 $X_1 \perp\!\!\!\!\not\perp X_3$ correlation and allowing accurate estimation of the direct $X_1 \leftrightarrow X_2$ and $X_2 \leftrightarrow X_3$ couplings. This is a concrete example of the Rung-2 advantage: only interventions, not observations, can resolve the ambiguity created by synchronization.

\subsection{Phillips Curve: Full Results}
\label{app:phillips}

\begin{figure}[h]
\centering
\begin{tikzpicture}[
    node distance=1cm,
    econ/.style={rectangle, draw, rounded corners, minimum width=1.5cm, minimum height=0.75cm, font=\scriptsize, thick, fill=gray!15},
    target/.style={rectangle, draw, rounded corners, minimum width=1.6cm, minimum height=0.8cm, font=\scriptsize, thick, fill=blue!20},
    arrow/.style={-{Stealth[length=2mm]}, thick},
    time/.style={font=\tiny, gray}
]
\node[econ] (UN) at (0,1.2) {UNRATE$_t$};
\node[econ] (FF) at (0,0)   {FEDFUNDS$_t$};
\node[econ] (MI) at (0,-1.2){MICH$_t$};
\node[target] (CPI) at (3.5,0) {CPI$_{t+1}$};
\draw[arrow] (UN) -- (CPI);
\draw[arrow] (FF) -- (CPI);
\draw[arrow] (MI) -- (CPI);
\node[time] at (0,-2)   {time $t$};
\node[time] at (3.5,-2) {time $t+1$};
\draw[gray, dashed] (1.75,1.8) -- (1.75,-1.8);
\node[font=\tiny, align=left] at (6.8,0.9) {Regimes:};
\node[font=\tiny, align=left] at (7.2,0.4) {1970s stagflation};
\node[font=\tiny, align=left] at (7.2,0.0) {Volcker disinflation};
\node[font=\tiny, align=left] at (7.2,-0.4){Great Moderation};
\node[font=\tiny, align=left] at (7.2,-0.8){2008 crisis};
\end{tikzpicture}
\caption{Phillips curve causal structure. Variables at time $t$ jointly determine CPI at $t+1$. ACE performs \emph{active data subset selection}: choosing which historical regimes to include in the training set, treating each regime as a ``virtual intervention'' on economic conditions. This is a Rung-2 analogue using archival data.}
\label{fig:phillips-scm}
\end{figure}

Using Federal Reserve Economic Data \citep{fred2024} (FRED, 1960--2023), we model $\text{CPI}_{t+1} = f(\text{UNRATE}_t, \text{FEDFUNDS}_t, \text{MICH}_t)$. ACE performs \emph{active data subset selection}: choosing which historical periods to query for training data, treating each regime as a virtual intervention on macroeconomic conditions.

\begin{table}[h]
\centering
\caption{Phillips curve out-of-sample CPI prediction MSE (N=5 seeds, 50 episodes). Regime selection percentages are policy-averaged fractions of training episodes spent in each historical regime.}
\label{tab:phillips-results}
\small
\begin{tabular}{@{}lccc@{}}
\toprule
Selection strategy & Out-of-sample MSE & Std & Regime preferences \\
\midrule
\textbf{ACE}               & \textbf{0.31} & 0.08 & 1970s 38\%, Volcker 24\%, Great Recession 19\% \\
Random regime              & 0.52          & 0.12 & Uniform across regimes \\
Sequential (chronological) & 0.44          & 0.10 & By date order \\
\bottomrule
\end{tabular}
\end{table}

ACE consistently prioritizes high-volatility regimes (1970s stagflation, Volcker disinflation, Great Recession) that expose nonlinear inflation dynamics absent from low-volatility eras---a $40\%$ reduction in out-of-sample MSE versus random regime selection. This demonstrates that preference-based experimental design generalizes to retrospective data-selection settings beyond controlled interventions, though we note this application departs from strict Rung-2 causal intervention.

\subsection{Lookahead vs.\ DPO Ablation}
\label{app:lookahead}

The No-DPO ablation (Section~\ref{app:ablation-detail}) isolates the contribution of learned proposal generation from the lookahead mechanism. Result: $2.10\pm0.11$, matching Round-Robin and confirming that lookahead alone provides no benefit. Concretely: if $K=4$ random candidates are generated, the probability that \emph{any} of them targets the collider parents ($X_1$ or $X_2$) with a well-chosen value is $\approx (2/5)^{4} \approx 10\%$---too low for the select-best procedure to reliably discover the collider strategy. The DPO-trained LM shifts the generation distribution so that strategic candidates appear in $>90\%$ of steps.

\subsection{Strategic Behavior: Intervention Distribution}
\label{app:strategic}

The learned policy concentrates $99.8\%$ of interventions on $X_1$ and $X_2$ (the two parents of the collider $X_3$). This is consistent with the causal theoretic optimum: to identify the nonlinear mechanism $f_3(X_1, X_2) = 0.5 X_1 - X_2 + \sin(X_2)$, the policy must vary both $X_1$ and $X_2$ across their value ranges while fixing exogenous noise. Under-sampled nodes $X_4$ and $X_5$ (the disconnected quadratic chain) are addressed primarily through the diversity reward, which ensures $L_{X_5}$ eventually decreases despite low explicit attention. Cross-seed consistency is high: the standard deviation of per-node intervention fractions across the 5 ACE seeds is $<0.01$ for nodes $X_1$ and $X_2$, confirming that strategy discovery is a robust consequence of the training procedure rather than a seed-dependent artifact.

% =======================================================================
\section{Failure Modes and Limitations}
\label{app:failures}
% =======================================================================

\subsection{Outlier Seed Analysis}
\label{app:outlier}

Seed 789 (5-node benchmark) produced $L_{X_5} \approx 1.73$ versus $0.02$--$0.22$ for other seeds. Two hypotheses: (i) initialization sensitivity for the quadratic mechanism $X_5 = 0.2 X_4^2$, where small initial MLP weights near zero may not escape a local minimum with flat gradient; (ii) the diversity reward in early episodes drove the policy away from $X_4$ before $L_{X_5}$ had been adequately reduced. We report the median ($0.61$) and note the outlier explicitly; the 5-node result is therefore conservative relative to the mean.

Seed 42 (30-node benchmark) exhibited training instability with loss oscillating $2.7$--$3.5$ throughout 40 completed episodes. The 30-node system has a longer effective horizon (the policy must identify 13+ colliders) and the checkpoint-resume mechanism may have caused optimizer state discontinuities. Seeds 123 and 456 converged cleanly.

\subsection{Graph Misspecification: Theoretical Failure Modes}
\label{app:misspec-failures}

Since ACE assumes a known causal graph, the practical cost of graph errors partitions cleanly by error type:

\textbf{Extra edge} $V_i \to V_j$ (spurious). The student learns $f_j(\text{Pa}_j \cup \{V_i\})$ instead of $f_j(\text{Pa}_j)$. The extra input $V_i$ receives near-zero weight in the MLP if $V_i$ truly has no causal effect on $V_j$; recovery is possible with sufficient interventions on $V_i$.

\textbf{Missing edge} $V_i \to V_j$ (absent). The student's mechanism $\hat{f}_j$ lacks the true parent $V_i$ as input. A lower bound on $L_j$ is set by the partial variance $\text{Var}(V_j - \hat{f}_j(\text{Pa}_j \setminus \{V_i\}))$. This is non-zero and non-recoverable; the student cannot decrease $L_j$ below this bound regardless of intervention strategy.

\textbf{Reversed edge} $V_i \to V_j$ (direction error). The student treats $V_j$ as a parent of $V_i$, so $\hat{f}_i$ depends on $V_j$. Under interventions on $V_i$, the true $V_j$ varies (since $V_j$ depends on $V_i$), but the student's conditioning set treats it as a predictor of $V_i$. This creates a circular dependency in the learner's prediction that no intervention strategy can resolve; it is the most damaging error type. A full empirical quantification across all four misspecification types was initiated on CURC but did not complete within our compute budget (L40 GPU OOM from $3$GB policy checkpoints). This remains our highest-priority follow-up experiment.

\subsection{Computational Cost}
\label{app:compute}

Each ACE training step requires:
\begin{itemize}[leftmargin=1.2em,topsep=0.2em,itemsep=0.1em]
\item $K=4$ LM forward passes (candidate generation): $\sim$2--3 min on A100.
\item $K=4$ cloned-learner training cycles (lookahead evaluation): $\sim$10--12 min.
\item One DPO update on the LM: $\sim$3--5 min.
\item Observational training (every 3 steps): $\sim$5 min amortized.
\end{itemize}
Total: $\sim$22 min per step on a single A100. A full run (171 episodes $\times$ 25 steps/episode with early stopping) takes $\sim$6--8 h per seed. Mitigation strategies for future work: smaller policy LMs (e.g., Qwen2.5-0.5B), learned value-of-lookahead to reduce $K$ adaptively, batch evaluation across candidates.

A note on checkpoint memory: the 1.5B-parameter LM in FP32 with optimizer state requires $\sim$17 GB per checkpoint, which caused OOM on CURC's L40 GPUs (46 GB VRAM) when combined with the KV cache and cloned learner. We implemented FP16 checkpointing (reducing to $\sim$3 GB) and eliminated optimizer state from checkpoints; these changes are in the repository.

% =======================================================================
\section{Reproducibility}
\label{app:reproducibility}
% =======================================================================

\subsection{Code and Data}

All code is released at \url{https://github.com/[anonymized]/ACE}. The repository includes:
\begin{itemize}[leftmargin=1.2em,topsep=0.2em,itemsep=0.1em]
\item Full ACE implementation (\texttt{ace\_experiments.py}) with FP16 checkpoint-resume logic.
\item Baselines (\texttt{baselines.py}): Random, Round-Robin, Max-Variance, PPO, Bayesian OED.
\item Experiments: 5-node and 30-node SCMs, Duffing oscillators, Phillips curve.
\item CPU baseline runner for 30-node (\texttt{scripts/runners/run\_30node\_cpu\_baselines.py}).
\item SLURM job scripts for CURC (\texttt{jobs/curc\_*.sh}).
\end{itemize}

\subsection{Random Seeds and Seed Assignment}

\begin{itemize}[leftmargin=1.2em,topsep=0.2em,itemsep=0.1em]
\item 5-node ACE: seeds 42, 123, 456, 789, 1011 (original Lambda Cloud runs).
\item 5-node baselines (heuristic + PPO): seeds 141, 271, 314, 577, 618 (CURC runs).
\item 5-node Bayesian OED: seeds 42, 123, 456 (CURC runs, higher per-seed cost).
\item 30-node ACE: seeds 42, 123, 456 (CURC \texttt{aa100} partition).
\item 30-node baselines: seeds 42, 123, 456, 789, 1011 (local CPU runs).
\item Duffing oscillators: seeds 42, 123, 456, 789, 1011 (N=5 all methods).
\item Phillips Curve: seeds 42, 123, 456, 789, 1011 (N=5 all methods).
\end{itemize}

\subsection{Compute Resources}

Main experiments ran on: CU Boulder CURC \texttt{aa100} partition (NVIDIA A100 80GB) for ACE GPU runs; CURC \texttt{al40} partition (NVIDIA L40 46GB) for baselines; Lambda Cloud A10 (24GB) for exploratory runs. Total: approximately 480 GPU-hours (ACE) and 120 CPU-hours (baselines). The Phillips Curve experiment required no GPU.

\subsection{Hyperparameters Summary}
\label{app:hyper-summary}

\begin{table}[h]
\centering
\caption{Complete hyperparameter table for all reported experiments.}
\label{tab:hyper-summary}
\small
\begin{tabular}{@{}lll@{}}
\toprule
Component & Hyperparameter & Value \\
\midrule
Policy & Model & Qwen2.5-1.5B \\
Policy & Sampling temperature & 0.7 \\
Policy & Candidates per step ($K$) & 4 \\
Policy & Supervised pretrain steps & 100 \\
Policy & Pretrain re-run every & 25 episodes \\
Policy & Optimizer & Adam, lr $1 \times 10^{-5}$ \\
\midrule
Learner & Architecture & 2-layer MLP, 64 hidden, ReLU \\
Learner & Optimizer & Adam, lr $2 \times 10^{-3}$ \\
Learner & Per-step training epochs & 100 \\
Learner & Observational training every & 3 steps \\
Learner & Observational batch size & 200 \\
\midrule
Reward & $\alpha$ (node importance weight) & 0.1 (\texttt{cov\_bonus=60.0}) \\
Reward & $\gamma$ (diversity weight) & 0.05 \\
Reward & $\beta$ (DPO temperature) & 0.1 \\
\midrule
Training & Max episodes & 200 (early stopping at convergence) \\
Training & Reference policy update every & 25 episodes \\
Training & Per-node convergence patience & 10 episodes \\
Training & Checkpoint interval & 1 step (rolling) + 10 steps (milestone) \\
\bottomrule
\end{tabular}
\end{table}

% =======================================================================
\section{Extended Related Work}
\label{app:related-extended}
% =======================================================================

\textbf{ABCI (Toth et al., 2022).} Active Bayesian Causal Inference \citep{toth2022active} maintains a posterior over causal graphs using continuous latent representations and Gaussian process mechanisms. The acquisition function selects experiments maximally informative about a downstream causal query. Compared to ACE: ABCI does not assume known structure; requires explicit posterior inference over graphs (combinatorial in $|V|$); targets causal-query inference rather than mechanism estimation. We adapt ABCI's acquisition-function philosophy as our Bayesian OED baseline by holding structure fixed and approximating the posterior over mechanism parameters.

\textbf{CBED (Tigas et al., 2022).} \citep{tigas2022interventions} jointly selects intervention targets and values at scale via Bayesian causal discovery, with demonstrations on the DREAM gene regulatory network (up to 20 variables). CBED addresses structure discovery; ACE addresses mechanism estimation. The two are complementary: CBED's structure-inference output can serve as ACE's input.

\textbf{Causal reasoning with LLMs (Kıcıman et al., 2023).} \citet{kiciman2023causal} conduct a systematic evaluation of GPT-4 on causal discovery benchmarks, finding substantial capability on pairwise causal direction tasks and commonsense causal queries. Crucially, the LM's causal reasoning is driven by semantic associations in pretraining data, not explicit causal structure learning. ACE exploits this: the LM's semantic prior biases proposal generation toward causally plausible interventions even before DPO training, which then refines the policy toward the specific system.

\textbf{Reasoning as World-Model Planning (Hao et al., 2023).} RAP \citep{hao2023reasoning} frames language model reasoning as planning in a world model, where the LM serves simultaneously as the reasoner and the forward model. This framing is directly relevant to ACE: the LM policy in ACE implicitly models the causal system (world model) while generating candidate interventions (planner), and DPO updates refine this combined representation toward the goal of efficient mechanism estimation.

\textbf{In-Context Learning as Bayesian Inference (Xie et al., 2022).} \citet{xie2022explanation} show that in-context learning corresponds to implicit Bayesian inference: the LM's predictions given a context are approximately equivalent to Bayesian predictions under a prior induced by the pretraining distribution. For ACE, this means the per-step prompt (containing per-node losses and recent interventions) implicitly updates the LM's posterior over the system, providing a form of Bayesian adaptation without explicit posterior computation.

\textbf{Multi-fidelity Bayesian active causal discovery (Zhang et al., 2023).} \citet{zhang2023bayesian} extend Bayesian OED to multi-fidelity oracle settings using mutual-information acquisition. The multi-fidelity insight is orthogonal to DPO-based learning and could in principle be incorporated into ACE's lookahead step.

\textbf{Sample-efficient Bayesian causal learning (Zhou et al., 2024).} \citet{zhou2024sample} address sample efficiency via polynomial-time DAG sampling. Their efficient structure learning is complementary to ACE's mechanism estimation: the two could be composed as a two-stage pipeline.
